\section{Conclusion}
\label{sec:conclusion}
In this paper, comparative analysis has been done on four small LLMs, Qwen3-0.6B, Qwen3-1.7B, Llama-3.2-1B-Instruct and Gemma-3-1B-it. These LLMs were fine-tuned with LoRA for binary code smell detection on the MLCQ dataset. The main goal is to find evidence that small light models can perform well in detecting code smells without the need for high resources and computation.
The results of this study show that Qwen3-1.7B achieved the best overall present-class F1-score of 0.6021, ranking first on three out of four smell types (Blob, Data Class, and Feature Envy). Gemma-3-1B-it came in second overall and did the best on detecting Long Method code smells. These results show that small language models can detect code smells well enough to be useful even when they are trained on a single free-tier GPU.
The performance of these language models varied a lot depending on the type of code smell. Long Method was the easiest to detect, with a F1-score of 0.80 across all models. This is probably because Long Method code smells are easier to spot since they are usually very long. On the other hand, Feature Envy was the hardest to detect with an average F1-score of only 0.30 since detecting feature envy requires understanding how parts of the code work together which can be more challenging.
One of the limitations of this study is that it only compared four language models that were fine-tuned in the same way. Different fine-tuning methods, prompt formats, or other settings could have given different results. Moreover, the MLCQ dataset focused on Java code from open-source projects, which means the results may not generalize to other languages or codebases.
Future work could explore different LoRA configurations and hyperparameter settings, since these can affect model performance. It would also be worth evaluating the models on datasets that include other programming languages.
